\documentclass[11pt]{article}

\usepackage[margin=1in]{geometry}
\usepackage{array}
\usepackage{booktabs}
\usepackage{longtable}
\usepackage{tabularx}
\usepackage{xcolor}
\usepackage{hyperref}

\hypersetup{
  colorlinks=true,
  linkcolor=black,
  urlcolor=blue
}

\setlength{\parindent}{0pt}
\setlength{\parskip}{6pt}
\renewcommand{\arraystretch}{1.18}

\newcolumntype{Y}{>{\raggedright\arraybackslash}X}
\newcommand{\markcell}[2]{\textbf{#1/#2}}

\begin{document}

\begin{center}
  {\Large Geography Paper 2 HLSL 2023: Student Work Marking Report}\\[4pt]
  {\normalsize Marked against \texttt{Geography\_paper\_2\_\_HLSL\_markscheme-2023.pdf}}
\end{center}

\section*{Overall Result}

\begin{tabularx}{\textwidth}{@{}p{0.34\textwidth} c c Y@{}}
\toprule
Section & Marks awarded & Maximum & Comment \\
\midrule
Question 1: Changing population & 9 & 10 & Lost one mark on the 0--14 estimate; ageing-society management is now credited. \\
Question 2: Global climate & 10 & 10 & Accurate and clearly developed throughout. \\
Question 3: Resource consumption and security & 8 & 10 & Strong oil answers; environmental impact of waste was not clearly environmental. \\
Question 4: Climate-change infographic & 9 & 10 & Accurate data answers and mostly balanced evaluation; final appraisal wording needs tightening. \\
Question 6: WFE nexus essay & 6 & 10 & Relevant, but not comparative enough for the higher band. \\
\midrule
\textbf{Total} & \textbf{42} & \textbf{50} & Main targets: resource estimates, impact categories, clearer appraisal wording, and comparative essay judgement. \\
\bottomrule
\end{tabularx}

\section*{Question-by-Question Marks}

\begin{longtable}{@{}p{0.12\textwidth} p{0.12\textwidth} p{0.31\textwidth} p{0.37\textwidth}@{}}
\toprule
Question & Mark & Awarded for & Feedback \\
\midrule
\endfirsthead
\toprule
Question & Mark & Awarded for & Feedback \\
\midrule
\endhead
1(a)(i) & \markcell{1}{1} & Correctly identifies age group 65--69. & Accurate reading of the population pyramid. \\
1(a)(ii) & \markcell{0}{1} & The answer gives 5 million. & The accepted range is 14--18 million. Add the male and female bars for ages 0--4, 5--9, and 10--14. \\
1(b) Way 1 & \markcell{2}{2} & Identifies disease prevention/improved health in an ageing population and links this to reduced pressure on health and long-term care systems. & This is credited because the markscheme accepts improving elderly health to reduce health/hospital costs. \\
1(b) Way 2 & \markcell{2}{2} & Gives reducing pension/retirement benefits to discourage early retirement. & Valid pension-system management with development through keeping older people economically active for longer. \\
1(c) & \markcell{4}{4} & Uses Nigeria consistently and explains two valid physical factors: arid hot desert/lack of water in the north and fertile southern land supporting agriculture. & Full marks. Precision note: northern Nigeria is better described as the Sahel or semi-arid north rather than the Sahara Desert proper. \\
\midrule
2(a) & \markcell{2}{2} & Defines global dimming as reduced solar radiation reaching Earth due to aerosols reflecting sunlight back to space. & Complete definition. \\
2(b) & \markcell{4}{4} & Explains higher HIC car use/fuel combustion and lower renewable access in developing or middle-income contexts leading to fossil-fuel use. & Both ways link economic development to international variation in greenhouse-gas emissions. \\
2(c) & \markcell{4}{4} & Explains that wealthy people have infrastructure/wealth to absorb food-price shocks and can choose lower-risk places to live. & Both reasons are developed enough for full marks. \\
\midrule
3(a)(i) & \markcell{1}{1} & Gives 185 million oil-equivalent tonnes for Japan. & Within the accepted range. \\
3(a)(ii) & \markcell{1}{1} & Identifies China as a country consuming more oil than it produces. & Correct. \\
3(b)(i) & \markcell{2}{2} & Explains electric vehicles reducing reliance on fossil-fuel cars and gasoline. & Valid reason and development for decreased per-person oil consumption in HICs. \\
3(b)(ii) & \markcell{2}{2} & Explains middle-class growth increasing car ownership, with cars generally gasoline-based. & Valid reason and development for increased per-person oil consumption in MICs. \\
3(c) & \markcell{2}{4} & The economic impact is fully credited: recycling/remanufacturing creates jobs and may reduce poverty. & The environmental impact answer focuses on health complications for e-waste workers. To earn those two marks, state an environmental impact such as soil or water contamination by heavy metals and explain the pathway. \\
\midrule
4(a)(i) & \markcell{1}{1} & Gives wildfires, droughts and extreme temperatures. & Correct. \\
4(a)(ii) & \markcell{1}{1} & Gives 2.5 billion. & Correct modal value. \\
4(b) & \markcell{2}{2} & Links increased greenhouse gases and temperature rise to more floods, droughts and hurricanes. & Full credit; the answer gives a valid reason and explains the increase in disasters. \\
4(c) & \markcell{5}{6} & Uses evidence that impacts vary by region, by economic cost, and by population displacement, while also recognizing that impacts are globally felt. & Strong body of evidence and enough implicit appraisal for credit, but the explicit final sentence is muddled. A clearer judgement should state that impacts are global but not uniform. \\
\midrule
6 & \markcell{6}{10} & Defines the WFE nexus, explains climate-change links to water, food and energy, and uses Kiribati as a relevant case example. & Focused and relevant, but too narrow for the 7--8 band because it does not properly weigh climate change against other major influences on the nexus. \\
\bottomrule
\end{longtable}

\section*{Teaching Feedback}

\textbf{What is working well}
\begin{itemize}
  \item Most short answers use clear cause-effect chains rather than isolated facts.
  \item Case/place detail is improving: Nigeria and Kiribati are both used in relevant contexts.
  \item The student can usually identify the correct command-term demand and develop an explanation for the second mark.
  \item The six-mark infographic response attempts a final judgement, which is essential for the appraisal mark.
\end{itemize}

\textbf{Main improvement targets}
\begin{itemize}
  \item For estimate questions, use the graph or resource directly and stay within the likely visual range.
  \item Keep impact categories separate: environmental impacts need pollution, ecosystem, biodiversity or physical-environment pathways; health impacts are usually social.
  \item Keep place terminology precise: use ``Sahel/semi-arid north of Nigeria'' rather than ``Sahara Desert in Nigeria''.
  \item In essay questions, compare the named factor with at least two other plausible factors before reaching a judgement. In Question 6, the answer needed direct comparison with influences such as population growth, rising incomes, changing diets, technology, governance or trade.
  \item Use data with context and make the final appraisal sentence unambiguous: explain what a figure proves rather than only quoting it.
\end{itemize}

\section*{Suggested Next Practice}

\begin{itemize}
  \item Redo Question 1(a)(ii) by summing the three 0--14 age bands from both sides of the pyramid.
  \item Rewrite Question 3(c)'s environmental impact in two sentences using this structure: impact, process, consequence for the receiving country.
  \item Plan Question 6 again using three comparison columns: climate change, population/economic growth, and technology/governance. End with a one-sentence ranking of relative importance.
\end{itemize}

\section*{Model Improvements}

\textbf{Question 1(b) stronger structure:}
Improve disease prevention and elderly healthcare so fewer older people need expensive hospital or long-term care. Raise the pension age so older people remain in work for longer, increasing tax revenue and reducing the number of years the state must pay pensions.

\textbf{Question 3(c) stronger environmental impact:}
E-waste dumps can pollute soil and groundwater when heavy metals leach out of broken electronic components. This damages ecosystems and can contaminate water supplies in receiving countries.

\textbf{Question 6 stronger essay route:}
Climate change is highly important because drought, sea-level rise and extreme weather can disrupt water supply, crop yields and hydropower at the same time. However, the nexus is also strongly shaped by population growth, changing diets, urbanization, governance and technology. A high-scoring answer should compare these factors directly and conclude where climate change is most important, where it is secondary, and why.

\section*{Notes}

\begin{itemize}
  \item The supplied student-work images show answers for Section A Questions 1--3, Section B Question 4, and Section C Question 6. Question 5 was not attempted and was not marked.
  \item This report records marking judgements and teaching feedback only; it does not reproduce the exam paper or markscheme.
\end{itemize}

\end{document}
